

%\subsection{Winograd Data Creation}
To better identify gender bias in coreference resolution systems, we build a new dataset centered on people entities referred by their occupations from a vocabulary of 40 occupations gathered from the US Department of Labor, shown in Table~\ref{tab:gender_table}.\footnote{Labor Force Statistics from the Current Population Survey, 2017. https://www.bls.gov/cps/cpsaat11.htm} 
We use the associated occupation statistics to determine what constitutes gender stereotypical roles (e.g. 90\% of nurses are women in this survey). 
Entities referred by different occupations are paired and used to construct test case scenarios.
%All co-reference decisions should be resolved without gender.
Sentences are duplicated using male and female pronouns, and contain equal numbers of correct co-reference decisions for all occupations.
In total, the dataset contains 3,160 sentences, split equally for development and test, created by researchers  familiar with the project.
%Each sentence in the dataset contains a pair of occupations (X and Y), one is male biased, the other is female biased. It  also contains one gender pronoun. 
%The dataset contains two types of sentences containing gendered pronoun references to one of the male or female stereotyped professions. 
Sentences were created to follow two prototypical templates but annotators were encouraged to come up with scenarios where entities could be interacting in plausible ways. %\footnote{For example, nurses treating construction workers for injuries.} 
Templates were selected to be challenging and designed to cover cases requiring semantics and syntax separately.\footnote{We do not claim this set of templates is complete, but that they provide representative examples that, pratically, show bias in existing systems.} 

%We considered two types of prototypical test scenarios based on rough templates:

%\vspace{-5pt}
\paragraph{Type 1: [\texttt{entity1}] [interacts with]   [\texttt{entity2}] [conjunction] [pronoun] [circumstances].} Prototypical WinoCoRef style sentences, where co-reference decisions must be made using world knowledge about given circumstances (Figure~\ref{fig:teaser}; Type 1).
Such examples are challenging because they contain no syntactic cues.

%\vspace{-5pt}
\paragraph{Type 2: [\texttt{entity1}] [interacts with]  [\texttt{entity2}] and then [interacts with] [pronoun] for [circumstances].} These tests can be resolved using syntactic information and understanding of the pronoun (Figure~\ref{fig:teaser}; Type 2). 
We expect systems to do well on such cases because both semantic and syntactic cues help disambiguation.

\begin{table}[t]
    \small
    \centering
    \begin{tabular}{|l|r|l|r|}
    \hline
    Occupation & \% & Occupation & \% \\
    \hline

carpenter &2 &  editor  &52  \\ 
mechanician &  4 & designers &54  \\
construction worker & 4 &  accountant  &61  \\ 
laborer & 4  & auditor &   61  \\
driver  & 6 & writer  &63  \\  
sheriff &14 & baker &65  \\ 
mover & 18   & clerk &72  \\ 
developer &20 &   cashier & 73  \\  
farmer  &22 &   counselors  &73  \\  
guard &22 &    attendant &  76  \\
chief & 27   & teacher & 78  \\ 
janitor & 34    &  sewer &80  \\
lawyer  &35 &     librarian &84  \\
cook &  38   & assistant & 85  \\
physician&  38 &   cleaner & 89  \\ 
ceo &39 &    housekeeper &89  \\
analyst &41 &    nurse & 90  \\
manager &43 &    receptionist  &90  \\
supervisor & 44   &  hairdressers  &92  \\
salesperson&  48 &   secretary & 95  \\
\hline
    \end{tabular}
    \caption{\small Occupations statistics used in WinoBias dataset, organized by the percent of people in the occupation who are reported as female. When woman dominate profession, we call linking the noun phrase referring to the job with female and male pronoun as `pro-stereotypical`, and `anti-stereotypical`, respectively. Similarly, if the occupation is male dominated, linking the noun phrase with the male and female pronoun is called, `pro-stereotypical` and `anti-steretypical`, respectively.}
    \label{tab:gender_table}
\end{table}
%one follows the routine ``X verb Y pronoun disambiguation''.  The other one is of the form: ``X verb1 Y and verb2 pronoun''. The second type of sentence is syntactically unambiguous. Considering the following examples:\\
%Type1: \textit{[The developer]$_{e_1}$ was unable to communicate with [the writer]$_{e_2}$ because [he]$_{pro\_1}$ only understands code.} \\
%Type2: \textit{[The writer]$_{e_1}$ handed [the developer]$_{e_2}$ a coffee and then smiled at [him]$_{pro\_2}$.}

%All the sentences cannot be resolved by just looking at the gender or the pronoun.

%Besides the two types, we also split the dataset to stereotype and non\_stereotype parts. The first one contains the sentences that following the statistical stereotype in professions, for instance, it will use ``she'' to refer to  ``the writer'' and ``he'' to refer to ``the developer''. The non\_stereotype dataset instead will use ``he'' to stand for ``the writer'' and ``she'' to mention ``the developer''.

%We evaluate the three systems on this new dataset.
%\vspace{-5pt}
\paragraph{Evaluation}
To evaluate models, we split the data in two sections: one where correct co-reference decisions require linking a gendered pronoun to an occupation stereotypically associated with the gender of the pronoun and one that requires linking to the anti-stereotypical occupation.
%If a model is gender biased, it can .
We say that a model passes the WinoBias test if for both Type 1 and Type 2 examples, pro-stereotyped and anti-stereotyped co-reference decisions are made with the same accuracy.
%A limitation of this test is that it can be used to prove a system is biased but not that it is unbiased.


%The results  in Table~\ref{tab:cluster_gender_ratio} shows that all the three systems perform better on the stereotyped dataset than on the non\_stereotyped dataset for both two types which means all of them have the  gender bias issue. The rule-based Stanford system is the most gender biased among the three systems. And UW system is more gender biased than the Berkeley system. The Stanford and Berkeley systems both rely on the gender file collected by~\cite{Bergsma:06}, but the Berkeley system will also rely on other features which makes it less biased than the Stanford System. However the UW system is the end-to-end system. All the information is collected from the dataset itself, which makes it more sensitive to the dataset.